\setcounterpageref{ex}{0}
\TaoPhieuKiemTra{04}

\setcounter{ex}{0}
\setcounter{bt}{0}

\textbf{\large I. PHẦN TRẮC NGHIỆM (3,0 điểm)}

{\itshape Khoanh tròn vào chữ cái đặt trước câu trả lời đúng. Mỗi câu đúng được 0,5 điểm.}

\begin{ex} % Câu 1
Số liền sau của số $89$ là
\choice
{$88$}
{$91$}
{$99$}
{\True $90$}
\loigiai{
Số liền sau của $89$ là: $89 + 1 = 90$.\\
\textbf{Đáp án đúng: D}
}
\end{ex}

\begin{ex} % Câu 2
Cho phép tính $42 + 16 = 58$. Các số $42$ và $16$ được gọi là
\choice
{Số bị trừ}
{Tổng}
{\True Số hạng}
{Hiệu}
\loigiai{
Các thành phần cộng với nhau gọi là các Số hạng.\\
\textbf{Đáp án đúng: C}
}
\end{ex}

\begin{ex} % Câu 3
Một can chứa đầy $5\text{ }l$ dầu. Đổ thêm vào $3\text{ }l$ nữa thì có tất cả
\choice
{\True $8\text{ }l$}
{$7\text{ }l$}
{$9\text{ }l$}
{$2\text{ }l$}
\loigiai{
$5 + 3 = 8\text{ }l$.\\
\textbf{Đáp án đúng: A}
}
\end{ex}

\begin{ex} % Câu 4
Hình tam giác có tất cả bao nhiêu cạnh?
\choice
{$2$ cạnh}
{\True $3$ cạnh}
{$4$ cạnh}
{$5$ cạnh}
\loigiai{
Hình tam giác có $3$ cạnh và $3$ đỉnh.\\
\textbf{Đáp án đúng: B}
}
\end{ex}

\begin{ex} % Câu 5
Điền dấu thích hợp: $13 - 5 \dots 12 - 4$.
\choice
{$>$}
{$<$}
{\True $=$}
{Không so sánh được}
\loigiai{
$13 - 5 = 8$; $12 - 4 = 8$. Vì $8 = 8$ nên điền dấu $=$.\\
\textbf{Đáp án đúng: C}
}
\end{ex}

\begin{ex} % Câu 6
Tính: $13 - 7 + 8 = \dots$
\choice
{$13$}
{$15$}
{$12$}
{\True $14$}
\loigiai{
$13 - 7 = 6$; $6 + 8 = 14$.\\
\textbf{Đáp án đúng: D}
}
\end{ex}

\vspace{4mm}
\textbf{\large II. PHẦN TỰ LUẬN (7,0 điểm)}

{\itshape Thí sinh trình bày chi tiết bài làm vào giấy kiểm tra.}

\begin{bt}[1,0 điểm]
Tính nhẩm:
\begin{multicols}{2}
\begin{enumerate}[a)]
  \item $8 + 8$
  \item $7 + 9$
  \item $16 - 5$
  \item $18 - 6$
\end{enumerate}
\end{multicols}
\loigiai{
\begin{itemize}
  \item $8 + 8 = 16$
  \item $7 + 9 = 16$
  \item $16 - 5 = 11$
  \item $18 - 6 = 12$
\end{itemize}
}
\end{bt}

\begin{bt}[1,5 điểm]
Đặt tính rồi tính:
\begin{multicols}{3}
\begin{enumerate}[a)]
  \item $31 + 44$
  \item $83 - 33$
  \item $49 + 48$
\end{enumerate}
\end{multicols}
\loigiai{
\begin{itemize}
  \item $31 + 44 = 75$
  \item $83 - 33 = 50$
  \item $49 + 48 = 97$
\end{itemize}
}
\end{bt}

\begin{bt}[1,0 điểm]
Tính giá trị biểu thức kèm đơn vị đo:
\begin{enumerate}[a)]
  \item $20\text{ kg} + 7\text{ kg} - 10\text{ kg}$
  \item $29\text{ }l - 15\text{ }l + 6\text{ }l$
\end{enumerate}
\loigiai{
\begin{enumerate}[a)]
  \item $20\text{ kg} + 7\text{ kg} - 10\text{ kg} = 27\text{ kg} - 10\text{ kg} = 17\text{ kg}$.
  \item $29\text{ }l - 15\text{ }l + 6\text{ }l = 14\text{ }l + 6\text{ }l = 20\text{ }l$.
\end{enumerate}
}
\end{bt}

\begin{bt}[1,5 điểm]
Bạn Châu có $19$ ngôi sao giấy. Bạn Dương có ít hơn bạn Châu $8$ ngôi sao giấy. Hỏi bạn Dương có bao nhiêu ngôi sao giấy?
\loigiai{
Số ngôi sao giấy bạn Dương có là:
\[ 19 - 8 = 11 \text{ (ngôi sao giấy)} \]
\textbf{Đáp số:} $11$ ngôi sao giấy.
}
\end{bt}

\begin{bt}[1,0 điểm]
Một thùng chứa đầy $39\text{ }l$ dầu ăn. Buổi sáng cửa hàng đã bán ra $16\text{ }l$. Hỏi trong thùng còn lại bao nhiêu lít dầu ăn?
\loigiai{
Trong thùng còn lại số lít dầu ăn là:
\[ 39 - 16 = 23\text{ }l \]
\textbf{Đáp số:} $23\text{ }l$ dầu ăn.
}
\end{bt}

\begin{bt}[1,0 điểm]
Hình vẽ có đoạn thẳng $AB$. Nếu lấy thêm hai điểm $C$ và $D$ nằm giữa $A$ và $B$ thì có tất cả bao nhiêu đoạn thẳng?
\loigiai{
Các điểm thẳng hàng theo thứ tự là $A, C, D, B$.\\
Các đoạn thẳng tạo thành gồm:
- Đoạn đơn: $AC, CD, DB$ (3 đoạn).
- Đoạn đôi: $AD, CB$ (2 đoạn).
- Đoạn ba: $AB$ (1 đoạn).\\
Tổng cộng: $3 + 2 + 1 = 6$ đoạn thẳng.\\
\textbf{Đáp số:} $6$ đoạn thẳng.
}
\end{bt}
